\section{Parameter Tiap Algoritma}

\begin{longtable}{
    @{}
    >{\raggedright\arraybackslash}p{2.5cm}          % Algoritma
    >{\raggedright\arraybackslash}p{5cm}          % Parameter kunci
    >{\raggedright\arraybackslash}p{5cm}          % Catatan
    @{}
}
\caption{Parameter kunci CoDA-EDA dan enam baseline} \\
\toprule
Algoritma & Parameter kunci (nilai) & Catatan \\
\midrule
\endfirsthead
\caption{Parameter kunci CoDA-EDA dan enam baseline (lanjutan)} \\
Algoritma & Parameter kunci (nilai) & Catatan \\
\midrule
\endhead
\bottomrule
\endfoot

CoDA-EDA &
$NP = 100$; $p_0 = 3/D$; $\mathit{mi\_rows} = 1500$; $w_{\mathrm{fp}} = 1,0$; $w_{\mathrm{wsc}} = 0,1$ &
Pohon Chow-Liu order-1 disusun via algoritma Prim; struktur hanya dipicu drift \\

compact-GA &
$NP = 100$; $p_0 = 3/D$; laju $\pm 1/NP$ &
Univariat; representasi biner identik dengan CoDA-EDA \\

compact-PSO &
$NP = 100$; $w = 0,7$; $c_1 = c_2 = 1,5$; $p_0 = 3/D$ &
Posisi-kecepatan diadaptasi ke ruang kebijakan diskret \\

BOA &
$NP = 100$; $\mathit{cand} = 8$; $\mathit{top\_vars} = 400$ &
Jaringan Bayes via BIC; kandidat induk dibatasi 400 variabel paling informatif agar layak pada $D$ besar \\

iBOA &
$NP = 100$; $k_{\mathrm{order}} = 2$ &
Pohon dependensi inkremental (Chow-Liu); pembanding langsung H3 \\

ACO &
$NP = 100$ (semut); $\rho = 0,1$ (evaporasi feromon) &
Parameter mengikuti nilai baku publikasi asli \\

UBK &
$NP = 100$; $\alpha$ adaptif $= 0,8 \cdot \max(0,05; 1 - t/200)$; derau $= 0,3$ &
Strategi berburu berbasis populasi \\

\end{longtable}

Prior kejarangan $p_0 = 3/D$ bersifat sadar-dimensi (aturan ABAC nyata umumnya $2$--$5$ pasangan atribut-nilai), sehingga jumlah bit aktif di awal pencarian tetap konstan meskipun $D$ berbeda antar-dataset.